In this work, we examined the accuracy, representability, numerical stability, and finite-shot behaviour of spin-resolved QSE across a chemically diverse molecular benchmark. 
In the statevector limit, QSE generally recovers the ground state well, but the present one-body operator pool is less reliable for low-lying excited states, with representability degrading as the active space grows. 
Much of this degradation is associated with the limited representation of states with substantial multiconfigurational character, reflecting a limitation of the present one-body operator pool. 
The resulting representability errors vary substantially across molecules and active spaces, reflecting their sensitivity to the electronic character of the target states and the orbitals included.
Among the ansätze considered, we find that UCCSD provides a robust compromise between spectral accuracy, circuit complexity, and spin purity. 
Present results motivate the exploration of larger operator pools, including higher-order excitations, together with benchmarks in larger active spaces.

Analysis of the overlap spectrum shows that the numerical instability of QSE originates largely from weakly supported operator-generated directions. 
Although these small-eigenvalue modes are highly sensitive to perturbations, they can still contain important spectral information and therefore cannot be discarded trivially. 
Under finite-shot sampling, we find that the optimal regularisation depends on both shot count and active-space size. 
For the present one-body QSE construction and measurement protocol, maintaining approximately constant finite-shot accuracy leads to an empirical total-shot scaling of roughly $\mathcal{O}\!\left(n_{\mathrm{orb}}^{11}\right)$.

Finally, while a naive interpretation of these results might suggest a steep measurement overhead for direct matrix estimation, this picture is softened by additional structure in the problem. 
We find that the projected QSE matrices exhibit substantial internal structure, as well as strong correlations between $\mathbf{H}$ and $\mathbf{S}$. 
We therefore envisage that structure-aware classical post-processing and joint reconstruction strategies could be used to significantly improve performance and recover comparable spectral accuracy from fewer measurements. 
Overall, this work highlights the current limitations of QSE in its straightforward form, while also pointing to promising directions in which its practical efficiency and robustness may be substantially improved.
